\documentclass[11pt,a4paper]{article}
\usepackage[utf8]{inputenc}
\usepackage[T1]{fontenc}
\usepackage[spanish,es-minimal,es-noshorthands]{babel}
\AtBeginDocument{\renewcommand{\sin}{\operatorname{sen}}\renewcommand{\arcsin}{\operatorname{arcsen}}}
\usepackage{lmodern}
\usepackage[sfdefault,lining]{FiraSans}
\usepackage[a4paper,top=2.3cm,bottom=2.2cm,left=1.8cm,right=1.8cm,headheight=1.3cm,headsep=0.4cm,footskip=1cm]{geometry}
\usepackage{amsmath,amssymb}
\usepackage{icomma}
\usepackage{graphicx}
\usepackage[table]{xcolor}
\usepackage{tikz}
\usepackage{fontawesome5}
\usepackage{ccicons}
\usepackage{fancyhdr}
\usepackage{lastpage}
\usepackage{tasks}
\usepackage{needspace}
\usepackage[most]{tcolorbox}
\usepackage{array}
\usepackage[hidelinks]{hyperref}
\AtBeginDocument{\hypersetup{pdftitle={Ficha de ejercicios de la unidad 6: Introducción al álgebra. Ecuaciones de primer grado. 1º ESO},pdfsubject={Matemáticas 1º ESO},pdfkeywords={lenguaje algebraico, expresiones algebraicas, valor numérico, monomios, polinomios, propiedad distributiva, ecuaciones de primer grado, paréntesis, denominadores, problemas}}}

\definecolor{navy}{HTML}{1B2748}
\definecolor{accent}{HTML}{E8603A}
\definecolor{teal}{HTML}{1F9E90}
\definecolor{soft}{HTML}{EEF1F8}
\definecolor{muted}{HTML}{5D6685}

\setlength{\parindent}{0pt}
\setlength{\parskip}{3pt}
\settasks{label=\alph*), label-format=\bfseries\color{navy}, label-width=1.3em, item-indent=1.8em, column-sep=0.6em, after-item-skip=5pt, before-skip=3pt, after-skip=4pt}

% ---------- capçalera i peu ----------
\newcommand{\llicencia}{\href{https://creativecommons.org/licenses/by-nc-sa/4.0/deed.es}{\ccbyncsa}\ \ Licencia Creative Commons Atribución-NoComercial-CompartirIgual 4.0 Internacional (CC BY-NC-SA 4.0)}
\pagestyle{fancy}
\fancyhf{}
\fancyhead[L]{\raisebox{-0.35\height}{\includegraphics[height=0.95cm]{logo-gva.png}}}
\fancyhead[C]{\small\color{muted}Ficha de ejercicios · Unidad 6 · 1\textsuperscript{o} ESO}
\fancyhead[R]{\raisebox{-0.35\height}{\includegraphics[height=0.95cm]{logo-vermellar.png}}}
\fancyfoot[L]{\scriptsize\color{muted}\llicencia}
\fancyfoot[R]{\small\color{muted}\thepage}
\renewcommand{\headrulewidth}{0.4pt}
\renewcommand{\headrule}{\color{muted!40}\hrule height \headrulewidth}
\fancypagestyle{primera}{\fancyhf{}\renewcommand{\headrulewidth}{0pt}
  \fancyfoot[L]{\scriptsize\color{muted}\llicencia}
  \fancyfoot[R]{\small\color{muted}\thepage}}

% ---------- nivells i exercicis ----------
\newcommand{\estrelles}[1]{{\color{accent}\ifcase#1\or\faStar\faStar[regular]\faStar[regular]\or\faStar\faStar\faStar[regular]\or\faStar\faStar\faStar\fi}}
\newcounter{exer}
\newcommand{\ex}[1]{\par\needspace{4\baselineskip}\medskip\stepcounter{exer}\noindent{\color{navy}\bfseries\theexer.}\ \ {\small\estrelles{#1}}\ \ }

\newcommand{\seccio}[2]{\par\needspace{10\baselineskip}\bigskip
  \begin{tcolorbox}[colback=navy,colframe=navy,coltext=white,arc=2mm,boxrule=0pt,left=3mm,right=3mm,top=1.2mm,bottom=1.2mm]
  {\large\bfseries #1\quad #2}
  \end{tcolorbox}\vspace{-2pt}}

\newcolumntype{C}[1]{>{\centering\arraybackslash}p{#1}}

% quadre de recordatori (només a la fitxa d'exercicis)
\newcommand{\recorda}[1]{\par\smallskip\begin{tcolorbox}[colback=accent!7,colframe=accent!55,arc=2mm,boxrule=0.5pt,left=2.5mm,right=2.5mm,top=1mm,bottom=1mm]\small{\color{accent}\faLightbulb[regular]}\ \textbf{Recuerda:}\ #1\end{tcolorbox}}
\newcommand{\res}[1]{#1}\newcommand{\rest}[1]{#1}

% =====================================================================
% Capçalera FentMates de les fitxes: un color per a cada curs.
% Els colors són els de la pàgina inicial (LEVELS a assets/js/materials.js);
% el logo ve de _eines/logo/pdf/ (python3 _eines/logo/genera.py).
%
% Ús, al final del preàmbul de la unitat:
%   % =====================================================================
% Capçalera FentMates de les fitxes: un color per a cada curs.
% Els colors són els de la pàgina inicial (LEVELS a assets/js/materials.js);
% el logo ve de _eines/logo/pdf/ (python3 _eines/logo/genera.py).
%
% Ús, al final del preàmbul de la unitat:
%   % =====================================================================
% Capçalera FentMates de les fitxes: un color per a cada curs.
% Els colors són els de la pàgina inicial (LEVELS a assets/js/materials.js);
% el logo ve de _eines/logo/pdf/ (python3 _eines/logo/genera.py).
%
% Ús, al final del preàmbul de la unitat:
%   \input{../comu/fentmates.tex}
%   \fmcurs{1eso}
% i, al començament del document:
%   \capcalera{Matemàtiques}{1r ESO}{Fitxa d'exercicis · Unitat 2}{Divisibilitat}
% Opcional abans de \input: \def\fmneed{14} (espai mínim abans de cada secció).
% =====================================================================
\makeatletter
\def\fm@logodir{../../logo/pdf}
\providecommand{\fmneed}{10}
\def\fm@cursnom{}

% clau del curs, color i color del text damunt del color
\newcommand{\fm@curs}[3]{\@namedef{fm@color@#1}{#2}\@namedef{fm@text@#1}{#3}}
\fm@curs{1eso}{E6B84A}{navy}
\fm@curs{2eso}{3B82F6}{white}
\fm@curs{3eso}{FF7A50}{white}
\fm@curs{4eso-a}{A78BFA}{white}
\fm@curs{4eso-b}{6D28D9}{white}
\fm@curs{1batx-mi}{9B1C1C}{white}
\fm@curs{1batx-ccss}{EF6B6B}{white}
\fm@curs{2batx-mii}{1E6B47}{white}
\fm@curs{2batx-ccss}{2BB3A3}{white}

\newcommand{\fmcurs}[1]{%
  \@ifundefined{fm@color@#1}{\PackageError{fentmates}{Curs desconegut: #1}{}}{}%
  \edef\fm@tmp{\noexpand\definecolor{curs}{HTML}{\@nameuse{fm@color@#1}}%
               \noexpand\colorlet{curstext}{\@nameuse{fm@text@#1}}}\fm@tmp
  \colorlet{cursfosc}{curs!72!black}%
  \def\fm@logo{\fm@logodir/vertical-#1.pdf}%
  \def\fm@marca{\fm@logodir/marca-#1.pdf}}

% ---------- capçalera de la primera pàgina ----------
% #1 matèria (petit), #2 curs, #3 tipus de document i unitat, #4 títol
\newcommand{\capcalera}[4]{%
  \gdef\fm@cursnom{#2}%
  \par\noindent
  \begin{tcolorbox}[colback=curs,colframe=curs,boxrule=0pt,arc=2.2mm,width=3.5cm,height=2.15cm,
      valign=center,left=3mm,right=2mm,top=1mm,bottom=1mm,nobeforeafter,box align=center]
    \color{curstext}{\footnotesize #1}\par\vspace{1pt}{\fontsize{19}{21}\selectfont\bfseries #2\par}%
  \end{tcolorbox}\hspace{5mm}%
  \begin{minipage}[c]{\dimexpr\linewidth-3.5cm-5mm-2.9cm\relax}
    {\color{cursfosc}\small\bfseries\MakeUppercase{#3}}\par\vspace{2pt}
    {\color{navy}\fontsize{20}{23}\selectfont\bfseries #4\par}
  \end{minipage}\hfill
  \begin{minipage}[c]{2.6cm}\raggedleft\includegraphics[width=2.6cm]{\fm@logo}\end{minipage}\par
  \vspace{2pt}}

% ---------- barres de secció del color del curs ----------
\renewcommand{\seccio}[2]{\par\needspace{\fmneed\baselineskip}\bigskip
  \begin{tcolorbox}[colback=curs,colframe=curs,coltext=curstext,arc=2mm,boxrule=0pt,left=3mm,right=3mm,top=1.2mm,bottom=1.2mm]
  {\large\bfseries #1\quad #2}
  \end{tcolorbox}\vspace{-2pt}}

% ---------- estrelles de dificultat del color del curs (a Batxillerat no se'n posen) ----------
\@ifundefined{estrelles}{}{%
  \renewcommand{\estrelles}[1]{{\color{curs}\ifcase#1\or\faStar\faStar[regular]\faStar[regular]\or\faStar\faStar\faStar[regular]\or\faStar\faStar\faStar\fi}}}

% ---------- capçalera i peu de les pàgines ----------
% primera pàgina: logos institucionals al peu
\fancypagestyle{primera}{\fancyhf{}\renewcommand{\headrulewidth}{0pt}
  \fancyfoot[L]{\raisebox{-0.32\height}{\includegraphics[height=0.78cm]{logo-gva.png}}\hspace{5mm}%
    \raisebox{-0.32\height}{\includegraphics[height=0.78cm]{logo-vermellar.png}}}
  \fancyfoot[R]{\scriptsize\color{muted}\href{https://creativecommons.org/licenses/by-nc-sa/4.0/deed.ca}{\ccbyncsa}\ \ CC BY-NC-SA 4.0 · fentmates.org\qquad\small\thepage}}
% la resta: el símbol a l'esquerra i el curs a la dreta, amb el títol al centre com abans
\fancyhead[L]{\raisebox{-0.3\height}{\includegraphics[height=0.62cm]{\fm@marca}}}
\fancyhead[R]{\small\bfseries\color{cursfosc}\fm@cursnom}
\renewcommand{\headrule}{\color{curs!45}\hrule height \headrulewidth}
\makeatother

%   \fmcurs{1eso}
% i, al començament del document:
%   \capcalera{Matemàtiques}{1r ESO}{Fitxa d'exercicis · Unitat 2}{Divisibilitat}
% Opcional abans de \input: \def\fmneed{14} (espai mínim abans de cada secció).
% =====================================================================
\makeatletter
\def\fm@logodir{../../logo/pdf}
\providecommand{\fmneed}{10}
\def\fm@cursnom{}

% clau del curs, color i color del text damunt del color
\newcommand{\fm@curs}[3]{\@namedef{fm@color@#1}{#2}\@namedef{fm@text@#1}{#3}}
\fm@curs{1eso}{E6B84A}{navy}
\fm@curs{2eso}{3B82F6}{white}
\fm@curs{3eso}{FF7A50}{white}
\fm@curs{4eso-a}{A78BFA}{white}
\fm@curs{4eso-b}{6D28D9}{white}
\fm@curs{1batx-mi}{9B1C1C}{white}
\fm@curs{1batx-ccss}{EF6B6B}{white}
\fm@curs{2batx-mii}{1E6B47}{white}
\fm@curs{2batx-ccss}{2BB3A3}{white}

\newcommand{\fmcurs}[1]{%
  \@ifundefined{fm@color@#1}{\PackageError{fentmates}{Curs desconegut: #1}{}}{}%
  \edef\fm@tmp{\noexpand\definecolor{curs}{HTML}{\@nameuse{fm@color@#1}}%
               \noexpand\colorlet{curstext}{\@nameuse{fm@text@#1}}}\fm@tmp
  \colorlet{cursfosc}{curs!72!black}%
  \def\fm@logo{\fm@logodir/vertical-#1.pdf}%
  \def\fm@marca{\fm@logodir/marca-#1.pdf}}

% ---------- capçalera de la primera pàgina ----------
% #1 matèria (petit), #2 curs, #3 tipus de document i unitat, #4 títol
\newcommand{\capcalera}[4]{%
  \gdef\fm@cursnom{#2}%
  \par\noindent
  \begin{tcolorbox}[colback=curs,colframe=curs,boxrule=0pt,arc=2.2mm,width=3.5cm,height=2.15cm,
      valign=center,left=3mm,right=2mm,top=1mm,bottom=1mm,nobeforeafter,box align=center]
    \color{curstext}{\footnotesize #1}\par\vspace{1pt}{\fontsize{19}{21}\selectfont\bfseries #2\par}%
  \end{tcolorbox}\hspace{5mm}%
  \begin{minipage}[c]{\dimexpr\linewidth-3.5cm-5mm-2.9cm\relax}
    {\color{cursfosc}\small\bfseries\MakeUppercase{#3}}\par\vspace{2pt}
    {\color{navy}\fontsize{20}{23}\selectfont\bfseries #4\par}
  \end{minipage}\hfill
  \begin{minipage}[c]{2.6cm}\raggedleft\includegraphics[width=2.6cm]{\fm@logo}\end{minipage}\par
  \vspace{2pt}}

% ---------- barres de secció del color del curs ----------
\renewcommand{\seccio}[2]{\par\needspace{\fmneed\baselineskip}\bigskip
  \begin{tcolorbox}[colback=curs,colframe=curs,coltext=curstext,arc=2mm,boxrule=0pt,left=3mm,right=3mm,top=1.2mm,bottom=1.2mm]
  {\large\bfseries #1\quad #2}
  \end{tcolorbox}\vspace{-2pt}}

% ---------- estrelles de dificultat del color del curs (a Batxillerat no se'n posen) ----------
\@ifundefined{estrelles}{}{%
  \renewcommand{\estrelles}[1]{{\color{curs}\ifcase#1\or\faStar\faStar[regular]\faStar[regular]\or\faStar\faStar\faStar[regular]\or\faStar\faStar\faStar\fi}}}

% ---------- capçalera i peu de les pàgines ----------
% primera pàgina: logos institucionals al peu
\fancypagestyle{primera}{\fancyhf{}\renewcommand{\headrulewidth}{0pt}
  \fancyfoot[L]{\raisebox{-0.32\height}{\includegraphics[height=0.78cm]{logo-gva.png}}\hspace{5mm}%
    \raisebox{-0.32\height}{\includegraphics[height=0.78cm]{logo-vermellar.png}}}
  \fancyfoot[R]{\scriptsize\color{muted}\href{https://creativecommons.org/licenses/by-nc-sa/4.0/deed.ca}{\ccbyncsa}\ \ CC BY-NC-SA 4.0 · fentmates.org\qquad\small\thepage}}
% la resta: el símbol a l'esquerra i el curs a la dreta, amb el títol al centre com abans
\fancyhead[L]{\raisebox{-0.3\height}{\includegraphics[height=0.62cm]{\fm@marca}}}
\fancyhead[R]{\small\bfseries\color{cursfosc}\fm@cursnom}
\renewcommand{\headrule}{\color{curs!45}\hrule height \headrulewidth}
\makeatother

%   \fmcurs{1eso}
% i, al començament del document:
%   \capcalera{Matemàtiques}{1r ESO}{Fitxa d'exercicis · Unitat 2}{Divisibilitat}
% Opcional abans de \input: \def\fmneed{14} (espai mínim abans de cada secció).
% =====================================================================
\makeatletter
\def\fm@logodir{../../logo/pdf}
\providecommand{\fmneed}{10}
\def\fm@cursnom{}

% clau del curs, color i color del text damunt del color
\newcommand{\fm@curs}[3]{\@namedef{fm@color@#1}{#2}\@namedef{fm@text@#1}{#3}}
\fm@curs{1eso}{E6B84A}{navy}
\fm@curs{2eso}{3B82F6}{white}
\fm@curs{3eso}{FF7A50}{white}
\fm@curs{4eso-a}{A78BFA}{white}
\fm@curs{4eso-b}{6D28D9}{white}
\fm@curs{1batx-mi}{9B1C1C}{white}
\fm@curs{1batx-ccss}{EF6B6B}{white}
\fm@curs{2batx-mii}{1E6B47}{white}
\fm@curs{2batx-ccss}{2BB3A3}{white}

\newcommand{\fmcurs}[1]{%
  \@ifundefined{fm@color@#1}{\PackageError{fentmates}{Curs desconegut: #1}{}}{}%
  \edef\fm@tmp{\noexpand\definecolor{curs}{HTML}{\@nameuse{fm@color@#1}}%
               \noexpand\colorlet{curstext}{\@nameuse{fm@text@#1}}}\fm@tmp
  \colorlet{cursfosc}{curs!72!black}%
  \def\fm@logo{\fm@logodir/vertical-#1.pdf}%
  \def\fm@marca{\fm@logodir/marca-#1.pdf}}

% ---------- capçalera de la primera pàgina ----------
% #1 matèria (petit), #2 curs, #3 tipus de document i unitat, #4 títol
\newcommand{\capcalera}[4]{%
  \gdef\fm@cursnom{#2}%
  \par\noindent
  \begin{tcolorbox}[colback=curs,colframe=curs,boxrule=0pt,arc=2.2mm,width=3.5cm,height=2.15cm,
      valign=center,left=3mm,right=2mm,top=1mm,bottom=1mm,nobeforeafter,box align=center]
    \color{curstext}{\footnotesize #1}\par\vspace{1pt}{\fontsize{19}{21}\selectfont\bfseries #2\par}%
  \end{tcolorbox}\hspace{5mm}%
  \begin{minipage}[c]{\dimexpr\linewidth-3.5cm-5mm-2.9cm\relax}
    {\color{cursfosc}\small\bfseries\MakeUppercase{#3}}\par\vspace{2pt}
    {\color{navy}\fontsize{20}{23}\selectfont\bfseries #4\par}
  \end{minipage}\hfill
  \begin{minipage}[c]{2.6cm}\raggedleft\includegraphics[width=2.6cm]{\fm@logo}\end{minipage}\par
  \vspace{2pt}}

% ---------- barres de secció del color del curs ----------
\renewcommand{\seccio}[2]{\par\needspace{\fmneed\baselineskip}\bigskip
  \begin{tcolorbox}[colback=curs,colframe=curs,coltext=curstext,arc=2mm,boxrule=0pt,left=3mm,right=3mm,top=1.2mm,bottom=1.2mm]
  {\large\bfseries #1\quad #2}
  \end{tcolorbox}\vspace{-2pt}}

% ---------- estrelles de dificultat del color del curs (a Batxillerat no se'n posen) ----------
\@ifundefined{estrelles}{}{%
  \renewcommand{\estrelles}[1]{{\color{curs}\ifcase#1\or\faStar\faStar[regular]\faStar[regular]\or\faStar\faStar\faStar[regular]\or\faStar\faStar\faStar\fi}}}

% ---------- capçalera i peu de les pàgines ----------
% primera pàgina: logos institucionals al peu
\fancypagestyle{primera}{\fancyhf{}\renewcommand{\headrulewidth}{0pt}
  \fancyfoot[L]{\raisebox{-0.32\height}{\includegraphics[height=0.78cm]{logo-gva.png}}\hspace{5mm}%
    \raisebox{-0.32\height}{\includegraphics[height=0.78cm]{logo-vermellar.png}}}
  \fancyfoot[R]{\scriptsize\color{muted}\href{https://creativecommons.org/licenses/by-nc-sa/4.0/deed.ca}{\ccbyncsa}\ \ CC BY-NC-SA 4.0 · fentmates.org\qquad\small\thepage}}
% la resta: el símbol a l'esquerra i el curs a la dreta, amb el títol al centre com abans
\fancyhead[L]{\raisebox{-0.3\height}{\includegraphics[height=0.62cm]{\fm@marca}}}
\fancyhead[R]{\small\bfseries\color{cursfosc}\fm@cursnom}
\renewcommand{\headrule}{\color{curs!45}\hrule height \headrulewidth}
\makeatother

\fmcurs{1eso}

\begin{document}
\thispagestyle{primera}

\capcalera{Matemáticas}{1\textsuperscript{o} ESO}{Ficha de ejercicios · Unidad 6}{Introducción al álgebra.\\ Ecuaciones de primer grado}
\vspace{10pt}
Nombre y apellidos: \rule{6.6cm}{0.4pt}\hfill Grupo: \rule{1.6cm}{0.4pt}\hfill Fecha: \rule{2.2cm}{0.4pt}
\vspace{8pt}

\begin{tcolorbox}[colback=soft,colframe=soft,arc=2mm,left=3mm,right=3mm,top=1.5mm,bottom=1.5mm]
\small Ejercicios graduados por niveles de dificultad:\quad \estrelles{1}\ básico \qquad \estrelles{2}\ medio \qquad \estrelles{3}\ avanzado.\\[2pt] Resuélvelos en el cuaderno mostrando todos los pasos y comprueba las soluciones de las ecuaciones. Haz los cálculos sin calculadora.
\end{tcolorbox}

% =====================================================
\seccio{1}{Lenguaje algebraico}

\ex{1} Expresa en lenguaje algebraico (usa la letra $x$):\begin{tasks}(2)
  \task El doble de un número
  \task Un número más 7
  \task El triple de un número menos 2
  \task La mitad de un número
  \task El cuadrado de un número
  \task Un número menos su cuarta parte
  \task El consecutivo de un número
  \task El doble de un número más su triple
\end{tasks}\recorda{Entre un número y una letra no escribimos el punto de multiplicar: $3\cdot x=3x$. «El doble de un número más 3» es $2x+3$, pero «el doble de la suma de un número y 3» es $2(x+3)$.}

\ex{1} Escribe con palabras qué expresa cada expresión algebraica:\begin{tasks}(3)
  \task $x+5$
  \task $3x$
  \task $x-1$
  \task $\dfrac{x}{3}$
  \task $2(x+4)$
  \task $x^2-1$
\end{tasks}

\ex{2} La edad de Nuria es $x$ años. Expresa en lenguaje algebraico:\begin{tasks}(2)
  \task La edad de Nuria dentro de 6 años.
  \task La edad de Nuria hace 4 años.
  \task La edad de su abuela, que tiene 5 veces la de Nuria.
  \task La edad de su primo, que tiene la mitad de años que ella.
  \task La suma de las edades de Nuria y su abuela.
  \task La edad de la abuela dentro de 10 años.
\end{tasks}

\ex{2} Expresa en lenguaje algebraico:\begin{tasks}(2)
  \task Dos números consecutivos.
  \task Un número par y el par siguiente.
  \task Un número impar.
  \task La suma de tres números consecutivos.
  \task El precio de $x$ cuadernos si cada uno cuesta $1,50$\,€.
  \task Los céntimos de $x$ monedas de 20 céntimos.
\end{tasks}

\ex{2} Escribe una fórmula para:\begin{tasks}(1)
  \task Los euros que pagamos por $n$ entradas de cine a 8\,€ cada una.
  \task El número de ruedas que hay en $c$ coches y $m$ motos.
  \task Los minutos que hay en $h$ horas.
  \task El perímetro y el área de un cuadrado de lado $l$.
\end{tasks}

\ex{3} Observa la serie de figuras hechas con puntos:\par\smallskip{\centering \begin{tikzpicture}[baseline=0]\fill[accent] (0,0.2) circle (0.09); \fill[accent] (0,0.6) circle (0.09);\fill[navy] (0.40,0.00) circle (0.09);\fill[navy] (0.40,0.40) circle (0.09);\fill[navy] (0.40,0.80) circle (0.09);\end{tikzpicture}\qquad\qquad \begin{tikzpicture}[baseline=0]\fill[accent] (0,0.2) circle (0.09); \fill[accent] (0,0.6) circle (0.09);\fill[navy] (0.40,0.00) circle (0.09);\fill[navy] (0.40,0.40) circle (0.09);\fill[navy] (0.40,0.80) circle (0.09);\fill[navy] (0.80,0.00) circle (0.09);\fill[navy] (0.80,0.40) circle (0.09);\fill[navy] (0.80,0.80) circle (0.09);\end{tikzpicture}\qquad\qquad \begin{tikzpicture}[baseline=0]\fill[accent] (0,0.2) circle (0.09); \fill[accent] (0,0.6) circle (0.09);\fill[navy] (0.40,0.00) circle (0.09);\fill[navy] (0.40,0.40) circle (0.09);\fill[navy] (0.40,0.80) circle (0.09);\fill[navy] (0.80,0.00) circle (0.09);\fill[navy] (0.80,0.40) circle (0.09);\fill[navy] (0.80,0.80) circle (0.09);\fill[navy] (1.20,0.00) circle (0.09);\fill[navy] (1.20,0.40) circle (0.09);\fill[navy] (1.20,0.80) circle (0.09);\end{tikzpicture}\par}\smallskip a) ¿Cuántos puntos tiene cada figura? ¿Cuántos tendrá la figura 4? \ b) Escribe una expresión para los puntos de la figura $n$. \ c) ¿Cuántos puntos tendrá la figura 10? ¿Y la figura 50?

\ex{3} Traduce al lenguaje algebraico:\begin{tasks}(1)
  \task La diferencia entre el triple de un número y su mitad.
  \task El doble de la suma de un número y su consecutivo.
  \task La suma de los cuadrados de dos números $a$ y $b$.
  \task El cuadrado de la suma de dos números $a$ y $b$.
\end{tasks}


% =====================================================
\seccio{2}{Expresiones algebraicas. Valor numérico}

\ex{1} Calcula el valor numérico de $4x-3$ para: \quad a) $x=2$ \quad b) $x=0$ \quad c) $x=-1$ \quad d) $x=5$\recorda{Para calcular el valor numérico, sustituimos la letra por el número y hacemos las operaciones respetando la jerarquía. Los números negativos, siempre entre paréntesis: si $x=-3$, \ $x^2=(-3)^2=9$.}

\ex{1} Calcula el valor numérico de cada expresión para los valores indicados:\begin{tasks}(2)
  \task $x^{2}+1$ \ para \ $x=3$
  \task $2a+b$ \ para \ $a=4,\ b=-3$
  \task $10-3x$ \ para \ $x=-2$
  \task $\dfrac{x}{2}+5$ \ para \ $x=8$
  \task $3(x-1)$ \ para \ $x=4$
  \task $a\cdot b-a$ \ para \ $a=-2,\ b=5$
\end{tasks}

\ex{2} Copia y completa la tabla de valores:\par\smallskip{\centering\renewcommand{\arraystretch}{1.35}\begin{tabular}{|c|C{0.9cm}|C{0.9cm}|C{0.9cm}|C{0.9cm}|C{0.9cm}|C{0.9cm}||}\hline $x$ & $-2$ & $-1$ & $0$ & $1$ & $2$ & $3$\\\hline $2x+1$ &  &  &  &  &  & \\\hline $x^2-x$ &  &  &  &  &  & \\\hline\end{tabular}\par}

\ex{2} Usa las fórmulas para calcular:\begin{tasks}(1)
  \task El área de un triángulo, $A=\dfrac{b\cdot h}{2}$, si $b=10$\,cm y $h=7$\,cm.
  \task El perímetro de un rectángulo, $P=2b+2h$, si $b=12$\,cm y $h=4,5$\,cm.
  \task El volumen de un cubo, $V=a^3$, si $a=4$\,cm.
\end{tasks}

\ex{2} Una compañía de teléfono cobra 5\,€ al mes más $0,10$\,€ por cada minuto de llamada. Escribe la expresión de lo que pagamos en un mes si hablamos $m$ minutos. ¿Cuánto pagamos si hablamos 120 minutos? ¿Y si hablamos 45?

\ex{3} ¿Qué expresión da el valor de $y$ en cada tabla?\par\smallskip{\centering\renewcommand{\arraystretch}{1.3}\begin{tabular}{|c|C{0.8cm}|C{0.8cm}|C{0.8cm}|C{0.8cm}|}\hline $x$ & 1 & 2 & 3 & 4\\\hline $y$ & 2 & 6 & 10 & 14\\\hline\end{tabular}\qquad\begin{tabular}{|c|C{0.8cm}|C{0.8cm}|C{0.8cm}|C{0.8cm}|}\hline $x$ & 0 & 1 & 2 & 3\\\hline $y$ & 1 & 2 & 5 & 10\\\hline\end{tabular}\par}

\ex{3} Busca probando un valor de $x$ que haga que: \ a) $3x+2$ valga 20; \ b) $x^2$ valga 49; \ c) $2x+5$ y $x+9$ tengan el mismo valor.


% =====================================================
\seccio{3}{Monomios y polinomios}

\ex{1} Copia y completa la tabla:\par\smallskip{\centering\renewcommand{\arraystretch}{1.35}\begin{tabular}{|c|C{2.2cm}|C{2.2cm}|C{1.4cm}|}\hline Monomio & Coeficiente & Parte literal & Grado\\\hline $6x^3$ & & & \\\hline $-a^2b$ & & & \\\hline $xy$ & & & \\\hline $-7$ & & & \\\hline $\tfrac{3}{4}x^2y^2$ & & & \\\hline $-z$ & & & \\\hline\end{tabular}\par}\recorda{En el monomio $-5x^2y$: el coeficiente es $-5$, la parte literal es $x^2y$ y el grado es la suma de los exponentes, $2+1=3$. Si no hay número, el coeficiente es 1 (o $-1$).}

\ex{1} ¿Cuáles de estas expresiones son monomios?\begin{tasks}(4)
  \task $5x^2$
  \task $x+1$
  \task $-3ab$
  \task $\dfrac{4}{x}$
  \task $8$
  \task $x^2y^3$
  \task $2^x$
  \task $\dfrac{x}{5}$
\end{tasks}

\ex{2} Agrupa los monomios semejantes: \quad $4x,\ \ -2x^2,\ \ 3xy,\ \ x,\ \ 5x^2,\ \ -xy,\ \ 7,\ \ -3x$

\ex{2} Escribe: \ a) dos monomios semejantes a $-2x^3$; \ b) un monomio de grado 5 con coeficiente 4; \ c) un monomio de grado 0; \ d) un monomio de grado 2 con dos letras diferentes.

\ex{2} Indica el número de términos, el grado y el término independiente de cada polinomio:\begin{tasks}(2)
  \task $3x^2-5x+1$
  \task $7-x^3$
  \task $x^4+2x^2-x$
  \task $5x-8$
\end{tasks}

\ex{3} Escribe un monomio semejante a $6x^2y$ que: \ a) sumado con él dé 0; \ b) sumado con él dé $10x^2y$; \ c) restado de él dé $x^2y$.

\ex{3} Calcula el valor numérico: \ a) del monomio $-3x^2y$ para $x=-2$, $y=5$; \ b) del polinomio $x^3-2x+1$ para $x=-1$; \ c) del polinomio $2x^2-3x$ para $x=\frac{1}{2}$.


% =====================================================
\seccio{4}{Operaciones con monomios. Reducir expresiones}

\ex{1} Calcula:\begin{tasks}(4)
  \task $5x+3x$
  \task $9a-4a$
  \task $x+x+x$
  \task $2x-7x$
  \task $4x^{2}+x^{2}$
  \task $-3ab+8ab$
  \task $6y-y$
  \task $-x-x$
\end{tasks}\recorda{Solo podemos sumar o restar monomios semejantes: sumamos los coeficientes y dejamos la parte literal. $3x+5x=8x$, pero $3x+2$ no se puede reducir.}

\ex{1} Reduce:\begin{tasks}(2)
  \task $3x+2+5x-1$
  \task $7-2x+4-x$
  \task $x^{2}+3x-2x^{2}+x$
  \task $5a+2b-3a+b$
\end{tasks}

\ex{1} Multiplica:\begin{tasks}(3)
  \task $3\cdot 5x$
  \task $2x\cdot 4x$
  \task $(-3a)\cdot 2a$
  \task $x\cdot x^{2}$
  \task $5x\cdot (-2y)$
  \task $(-4x^{2})\cdot (-x)$
\end{tasks}

\ex{2} Aplica la propiedad distributiva:\begin{tasks}(3)
  \task $4(x+3)$
  \task $2(5x-1)$
  \task $-3(x+2)$
  \task $-(4x-7)$
  \task $x(x+5)$
  \task $2x(3x-1)$
\end{tasks}\recorda{$a(b+c)=ab+ac$: multiplicamos por \emph{cada} término del paréntesis. Un signo menos delante cambia todos los signos: $-(x-5)=-x+5$.}

\ex{2} Quita paréntesis y reduce:\begin{tasks}(2)
  \task $2(x+4)+3x$
  \task $5x-(2x+3)$
  \task $3(x-2)+2(x+5)$
  \task $4(2x+1)-(x-6)$
\end{tasks}

\ex{2} Halla el error en cada igualdad y corrígela:\begin{tasks}(2)
  \task $3x+4=7x$
  \task $2(x+5)=2x+5$
  \task $-(x-3)=-x-3$
  \task $x\cdot x=2x$
\end{tasks}

\ex{3} Reduce:\begin{tasks}(1)
  \task $3(2x-1)-2(x-4)+x$
  \task $5-2(3x+1)+4(x-2)$
  \task $x(x+2)-3(x-1)$
\end{tasks}

\ex{3} Escribe de forma reducida el perímetro de: \ a) un rectángulo de base $2x+1$ y altura $x-1$; \ b) un triángulo de lados $x$, $x+3$ y $2x-1$. Calcúlalos para $x=4$.

\ex{3} Adivinanza: piensa un número, súmale 5, multiplica el resultado por 2, réstale 10 y divide entre 2. ¿Qué obtienes? Explica por qué pasa siempre usando el álgebra.


% =====================================================
\seccio{5}{Ecuaciones. Elementos y solución}

\ex{1} Indica los miembros, los términos, la incógnita y el grado de cada ecuación: \quad a) $5x-3=2x+9$ \quad b) $4(a+1)=a-2$\recorda{Los \textbf{miembros} son las expresiones de cada lado del signo $=$; los \textbf{términos}, cada sumando; la \textbf{incógnita}, la letra; y la \textbf{solución}, el valor que hace cierta la igualdad.}

\ex{1} Comprueba si el valor indicado es solución de la ecuación:\begin{tasks}(2)
  \task $3x-1=8$, \ $x=3$
  \task $2x+5=x+1$, \ $x=-4$
  \task $5-x=2x$, \ $x=2$
  \task $\dfrac{x}{3}+1=x-3$, \ $x=6$
\end{tasks}

\ex{1} Halla mentalmente la solución:\begin{tasks}(3)
  \task $x+4=10$
  \task $x-3=8$
  \task $5x=30$
  \task $\dfrac{x}{2}=7$
  \task $12-x=5$
  \task $2x+1=11$
\end{tasks}

\ex{2} ¿Cuál de los valores $-2$, $0$, $1$ o $3$ es la solución de cada ecuación?\begin{tasks}(2)
  \task $4x-3=2x+3$
  \task $2(x+1)=x$
  \task $7x=x$
  \task $3-x=2x$
\end{tasks}

\ex{2} Di si cada igualdad es una identidad o una ecuación:\begin{tasks}(2)
  \task $2(x+1)=2x+2$
  \task $3x+1=3x+2$
  \task $x+x=2x$
  \task $4x-x=12$
\end{tasks}

\ex{3} Escribe una ecuación de primer grado: \ a) que tenga por solución $x=5$; \ b) que tenga por solución $x=-2$ y un paréntesis; \ c) que no tenga solución.


% =====================================================
\seccio{6}{Ecuaciones equivalentes. Resolución}

\ex{1} Resuelve:\begin{tasks}(4)
  \task $x+8=15$
  \task $x-6=-2$
  \task $x+9=4$
  \task $-5+x=-11$
\end{tasks}\recorda{Transposición: lo que \emph{suma} pasa \emph{restando} y lo que \emph{resta} pasa \emph{sumando}; lo que \emph{multiplica} pasa \emph{dividiendo} (con su signo) y lo que \emph{divide} pasa \emph{multiplicando}.}

\ex{1} Resuelve:\begin{tasks}(3)
  \task $6x=42$
  \task $-4x=20$
  \task $9x=-27$
  \task $\dfrac{x}{5}=3$
  \task $\dfrac{x}{4}=-3$
  \task $-x=8$
\end{tasks}

\ex{1} Resuelve:\begin{tasks}(4)
  \task $2x+3=11$
  \task $5x-4=21$
  \task $3x+10=1$
  \task $7-2x=13$
\end{tasks}

\ex{2} Resuelve y comprueba la solución:\begin{tasks}(2)
  \task $5x+2=3x+10$
  \task $8x-5=6x+1$
  \task $4x+7=x-8$
  \task $2x-9=5x+3$
\end{tasks}

\ex{2} Resuelve:\begin{tasks}(2)
  \task $3x+4-x=x+9$
  \task $6-2x+5=3x+1$
  \task $x+2x+3x=4x+10$
  \task $9x-3-4x=2x+12$
\end{tasks}

\ex{2} Resuelve (la solución puede ser una fracción: dala irreducible):\begin{tasks}(4)
  \task $4x=10$
  \task $6x+1=2x+3$
  \task $3x-2=5$
  \task $10x=-4$
\end{tasks}

\ex{3} Resuelve. ¿Qué pasa en las dos últimas?\begin{tasks}(2)
  \task $2x+5-7x=3-x-6$
  \task $12-3x+x=4x-6$
  \task $x+1=x+5$
  \task $3x+6-x=2x+6$
\end{tasks}

\ex{3} Escribe la ecuación que representa la balanza y halla cuánto pesa cada caja $x$ (las bolas pesan 1\,kg).\par\smallskip{\centering \begin{tikzpicture}[scale=0.85]\draw[fill=soft,draw=muted] (3,2.2)--(2.6,-0.4)--(3.4,-0.4)--cycle; \draw[line width=2pt] (0.6,2.2)--(5.4,2.2); \fill[accent] (3,2.2) circle (0.09);\draw[muted] (1.20,2.2)--(-0.10,0.5) (1.20,2.2)--(2.50,0.5);\draw[thick,fill=soft] (-0.10,0.5) arc (180:360:1.3 and 0.25);\draw[fill=navy!80,draw=navy] (0.05,0.53) rectangle +(0.38,0.38); \node[text=white,font=\small\itshape] at (0.24,0.72) {x};\draw[fill=navy!80,draw=navy] (0.53,0.53) rectangle +(0.38,0.38); \node[text=white,font=\small\itshape] at (0.72,0.72) {x};\draw[fill=navy!80,draw=navy] (1.01,0.53) rectangle +(0.38,0.38); \node[text=white,font=\small\itshape] at (1.20,0.72) {x};\draw[fill=accent!60,draw=accent] (1.68,0.72) circle (0.18); \node[font=\tiny] at (1.68,0.72) {1};\draw[fill=accent!60,draw=accent] (2.16,0.72) circle (0.18); \node[font=\tiny] at (2.16,0.72) {1};\draw[fill=accent!60,draw=accent] (0.24,1.18) circle (0.18); \node[font=\tiny] at (0.24,1.18) {1};\draw[fill=accent!60,draw=accent] (0.72,1.18) circle (0.18); \node[font=\tiny] at (0.72,1.18) {1};\draw[muted] (4.80,2.2)--(3.50,0.5) (4.80,2.2)--(6.10,0.5);\draw[thick,fill=soft] (3.50,0.5) arc (180:360:1.3 and 0.25);\draw[fill=navy!80,draw=navy] (3.65,0.53) rectangle +(0.38,0.38); \node[text=white,font=\small\itshape] at (3.84,0.72) {x};\draw[fill=accent!60,draw=accent] (4.32,0.72) circle (0.18); \node[font=\tiny] at (4.32,0.72) {1};\draw[fill=accent!60,draw=accent] (4.80,0.72) circle (0.18); \node[font=\tiny] at (4.80,0.72) {1};\draw[fill=accent!60,draw=accent] (5.28,0.72) circle (0.18); \node[font=\tiny] at (5.28,0.72) {1};\draw[fill=accent!60,draw=accent] (5.76,0.72) circle (0.18); \node[font=\tiny] at (5.76,0.72) {1};\draw[fill=accent!60,draw=accent] (3.84,1.18) circle (0.18); \node[font=\tiny] at (3.84,1.18) {1};\draw[fill=accent!60,draw=accent] (4.32,1.18) circle (0.18); \node[font=\tiny] at (4.32,1.18) {1};\draw[fill=accent!60,draw=accent] (4.80,1.18) circle (0.18); \node[font=\tiny] at (4.80,1.18) {1};\draw[fill=accent!60,draw=accent] (5.28,1.18) circle (0.18); \node[font=\tiny] at (5.28,1.18) {1};\draw[fill=accent!60,draw=accent] (5.76,1.18) circle (0.18); \node[font=\tiny] at (5.76,1.18) {1};\draw[fill=accent!60,draw=accent] (3.84,1.64) circle (0.18); \node[font=\tiny] at (3.84,1.64) {1};\end{tikzpicture}\par}


% =====================================================
\seccio{7}{Ecuaciones con paréntesis}

\ex{1} Resuelve:\begin{tasks}(2)
  \task $3(x+2)=18$
  \task $5(x-1)=20$
  \task $2(x+7)=x+12$
  \task $-2(x-3)=10$
\end{tasks}\recorda{Primero quitamos paréntesis (propiedad distributiva) y después resolvemos como siempre. Cuidado con el signo menos: $-2(x-3)=-2x+6$.}

\ex{2} Resuelve:\begin{tasks}(2)
  \task $4(x-1)=2(x+3)$
  \task $3(2x+1)=5x+7$
  \task $6-(x+2)=x$
  \task $2x-(5-x)=10$
\end{tasks}

\ex{2} Resuelve:\begin{tasks}(2)
  \task $5(x+2)-3x=16$
  \task $7x-2(x+4)=12$
  \task $3(x-4)+2(x+1)=5$
  \task $4(2x-3)=5(x+3)$
\end{tasks}

\ex{3} Resuelve. ¿Qué pasa en la última ecuación?\begin{tasks}(1)
  \task $2(3x-1)-3(x+2)=x+4$
  \task $5-2(x-3)=3(x-8)+5$
  \task $4(x+1)-3(x-2)=2(x+5)$
  \task $3(x-1)-(2x+5)=x-8$
\end{tasks}

\ex{3} Resuelve:\begin{tasks}(1)
  \task $x-2(x-5)=3(4-x)+6$
  \task $6(x+1)-4(2x-1)=0$
\end{tasks}


% =====================================================
\seccio{8}{Ecuaciones con denominadores}

\ex{1} Resuelve:\begin{tasks}(4)
  \task $\dfrac{x}{3}=5$
  \task $\dfrac{x}{4}+2=5$
  \task $\dfrac{2x}{3}=8$
  \task $\dfrac{x}{2}-1=4$
\end{tasks}\recorda{Multiplicamos \emph{todos} los términos de los dos miembros por el m.c.m. de los denominadores (también los que no tienen denominador). Si un numerador tiene más de un término, lo ponemos entre paréntesis.}

\ex{2} Resuelve:\begin{tasks}(4)
  \task $\dfrac{x}{2}+\dfrac{x}{4}=6$
  \task $\dfrac{x}{3}-\dfrac{x}{6}=2$
  \task $\dfrac{x}{2}+\dfrac{x}{3}=10$
  \task $x-\dfrac{x}{5}=8$
\end{tasks}

\ex{2} Resuelve:\begin{tasks}(4)
  \task $\dfrac{x+3}{2}=5$
  \task $\dfrac{x-1}{3}=x-5$
  \task $\dfrac{2x+1}{5}=3$
  \task $\dfrac{x}{2}=\dfrac{x+6}{5}$
\end{tasks}

\ex{2} Resuelve:\begin{tasks}(2)
  \task $\dfrac{x}{2}+3=\dfrac{x}{3}+5$
  \task $\dfrac{x+1}{4}=\dfrac{x-1}{2}$
\end{tasks}

\ex{3} Resuelve:\begin{tasks}(3)
  \task $\dfrac{x+2}{3}-\dfrac{x-1}{4}=2$
  \task $x-\dfrac{x+2}{3}=6$
  \task $\dfrac{3x-1}{2}-\dfrac{x+5}{4}=\dfrac{x}{4}$
\end{tasks}

\ex{3} Resuelve:\begin{tasks}(2)
  \task $\dfrac{x}{2}+\dfrac{x}{4}+\dfrac{x}{8}=14$
  \task $\dfrac{2(x-1)}{3}=\dfrac{x+4}{2}$
\end{tasks}


% =====================================================
\seccio{9}{Resolución de problemas con ecuaciones}

\ex{1} Si a un número le sumamos 12, obtenemos 30. ¿Cuál es el número?\recorda{Pasos: 1) lee y entiende el enunciado; 2) elige la incógnita ($x=\ldots$); 3) plantea la ecuación; 4) resuélvela; 5) comprueba la solución y responde a la pregunta.}

\ex{1} El triple de un número menos 8 es 25. ¿Cuál es el número?

\ex{1} La suma de dos números consecutivos es 73. ¿Cuáles son?

\ex{2} La suma de tres números consecutivos es 159. ¿Cuáles son?

\ex{2} Laura tiene 6 años más que su hermano. Entre los dos suman 30 años. ¿Cuántos años tiene cada uno?

\ex{2} Un bolígrafo y una libreta cuestan juntos $4,50$\,€. La libreta cuesta el doble que el bolígrafo. ¿Cuánto cuesta cada cosa?

\ex{2} El perímetro de un rectángulo es 46\,cm y la base mide 5\,cm más que la altura. ¿Cuánto miden los lados?

\ex{2} Repartimos 90 cromos entre tres amigos. El segundo recibe el doble que el primero y el tercero, el triple que el primero. ¿Cuántos recibe cada uno?

\ex{3} Una madre tiene 42 años y su hijo, 10. ¿Dentro de cuántos años la madre tendrá el triple de edad que el hijo?

\ex{3} La mitad de un número más su cuarta parte es 21. ¿Cuál es el número?

\ex{3} En una clase hay 28 alumnos. Si hubiera 4 chicas más, habría tantas chicas como chicos. ¿Cuántas chicas y cuántos chicos hay?

\ex{3} En un aparcamiento hay coches y motos: en total, 30 vehículos y 96 ruedas. ¿Cuántos coches y cuántas motos hay?

\ex{3} Los ángulos de un triángulo miden $x$, $2x$ y $3x$. ¿Cuánto mide cada ángulo? (Recuerda que los tres ángulos de un triángulo suman $180^\circ$.)

\vfill
\begin{center}\color{muted}\small\faCheckCircle[regular]\ \ ¡Ánimo! Si te cuesta algún ejercicio, repasa los apuntes interactivos de la unidad.\end{center}

\end{document}
